%%
%% Automatically generated file from DocOnce source
%% (https://github.com/hplgit/doconce/)
%%

% #define PREAMBLE

% #ifdef PREAMBLE
%-------------------- begin preamble ----------------------

\documentclass[%
oneside,                 % oneside: electronic viewing, twoside: printing
final,                   % draft: marks overfull hboxes, figures with paths
10pt,french]{article}

\listfiles               %  print all files needed to compile this document

\usepackage{relsize,makeidx,color,setspace,amsmath,amsfonts,amssymb}
\usepackage[table]{xcolor}
\usepackage{bm,ltablex,microtype}

\usepackage[pdftex]{graphicx}

% Packages for typesetting blocks of computer code
\usepackage{fancyvrb,framed,moreverb}

% Define colors
\definecolor{orange}{cmyk}{0,0.4,0.8,0.2}
\definecolor{tucorange}{rgb}{1.0,0.64,0}
\definecolor{darkorange}{rgb}{.71,0.21,0.01}
\definecolor{darkgreen}{rgb}{.12,.54,.11}
\definecolor{myteal}{rgb}{.26, .44, .56}
\definecolor{gray}{gray}{0.45}
\definecolor{mediumgray}{gray}{.8}
\definecolor{lightgray}{gray}{.95}
\definecolor{brown}{rgb}{0.54,0.27,0.07}
\definecolor{purple}{rgb}{0.5,0.0,0.5}
\definecolor{darkgray}{gray}{0.25}
\definecolor{darkblue}{rgb}{0,0.08,0.45}
\definecolor{darkblue2}{rgb}{0,0,0.8}
\definecolor{lightred}{rgb}{1.0,0.39,0.28}
\definecolor{lightgreen}{rgb}{0.48,0.99,0.0}
\definecolor{lightblue}{rgb}{0.53,0.81,0.92}
\definecolor{lightblue2}{rgb}{0.3,0.3,1.0}
\definecolor{lightpurple}{rgb}{0.87,0.63,0.87}
\definecolor{lightcyan}{rgb}{0.5,1.0,0.83}

\colorlet{comment_green}{green!50!black}
\colorlet{string_red}{red!60!black}
\colorlet{keyword_pink}{magenta!70!black}
\colorlet{indendifier_green}{green!70!white}

% Backgrounds for code
\definecolor{cbg_gray}{rgb}{.95, .95, .95}
\definecolor{bar_gray}{rgb}{.92, .92, .92}

\definecolor{cbg_yellowgray}{rgb}{.95, .95, .85}
\definecolor{bar_yellowgray}{rgb}{.95, .95, .65}

\colorlet{cbg_yellow2}{yellow!10}
\colorlet{bar_yellow2}{yellow!20}

\definecolor{cbg_yellow1}{rgb}{.98, .98, 0.8}
\definecolor{bar_yellow1}{rgb}{.98, .98, 0.4}

\definecolor{cbg_red1}{rgb}{1, 0.85, 0.85}
\definecolor{bar_red1}{rgb}{1, 0.75, 0.85}

\definecolor{cbg_blue1}{rgb}{0.87843, 0.95686, 1.0}
\definecolor{bar_blue1}{rgb}{0.7,     0.95686, 1}

%\setlength{\fboxsep}{-1.5mm}  % adjust cod_vpad/pro_vpad background box

%% Background for code blocks (parameter is color name)

%% pro/cod_vpad: gives some vertical padding before and after the text
%% (but has more simplistic code than _cod/pro_tight+cod/pro).
%% pro/cod_vpad can be used to enclose Verbatim or lst begin/end for code.
%% pro/cod calls _pro/cod_tight and has very little vertical padding,
%% used to enclose Verbatim and other begin/end for code.
%% (pro/cod is what the ptex2tex program could produce with the
%% Blue/BlueBar definitions in .ptex2tex.cfg.)

\newenvironment{cod_vpad}[1]{
   \def\FrameCommand{\colorbox{#1}}
   \MakeFramed{\FrameRestore}}
   {\endMakeFramed}

\newenvironment{_cod_tight}[1]{
   \def\FrameCommand{\colorbox{#1}}
   \FrameRule0.6pt\MakeFramed {\FrameRestore}\vskip3mm}
   {\vskip0mm\endMakeFramed}

\newenvironment{cod}[1]{
\bgroup\rmfamily
\fboxsep=0mm\relax
\begin{_cod_tight}{#1}
\list{}{\parsep=-2mm\parskip=0mm\topsep=0pt\leftmargin=2mm
\rightmargin=2\leftmargin\leftmargin=4pt\relax}
\item\relax}
{\endlist\end{_cod_tight}\egroup}

%% Background for complete program blocks (parameter 1 is color name
%% for background, parameter 2 is color for left bar)
\newenvironment{pro_vpad}[2]{
   \def\FrameCommand{\color{#2}\vrule width 1mm\normalcolor\colorbox{#1}}
   \MakeFramed{\FrameRestore}}
   {\endMakeFramed}

\newenvironment{_pro_tight}[2]{
   \def\FrameCommand{\color{#2}\vrule width 1mm\normalcolor\colorbox{#1}}
   \FrameRule0.6pt\MakeFramed {\advance\hsize-2mm\FrameRestore}\vskip3mm}
   {\vskip0mm\endMakeFramed}

\newenvironment{pro}[2]{
\bgroup\rmfamily
\fboxsep=0mm\relax
\begin{_pro_tight}{#1}{#2}
\list{}{\parsep=-2mm\parskip=0mm\topsep=0pt\leftmargin=2mm
\rightmargin=2\leftmargin\leftmargin=4pt\relax}
\item\relax}
{\endlist\end{_pro_tight}\egroup}

\usepackage{minted}
\usemintedstyle{default}

\usepackage[T1]{fontenc}
%\usepackage[latin1]{inputenc}
\usepackage{ucs}
\usepackage[utf8x]{inputenc}

\usepackage{lmodern}         % Latin Modern fonts derived from Computer Modern

% Hyperlinks in PDF:
\definecolor{linkcolor}{rgb}{0,0,0.4}
\usepackage{hyperref}
\hypersetup{
    breaklinks=true,
    colorlinks=true,
    linkcolor=linkcolor,
    urlcolor=linkcolor,
    citecolor=black,
    filecolor=black,
    %filecolor=blue,
    pdfmenubar=true,
    pdftoolbar=true,
    bookmarksdepth=3   % Uncomment (and tweak) for PDF bookmarks with more levels than the TOC
    }
%\hyperbaseurl{}   % hyperlinks are relative to this root

\setcounter{tocdepth}{2}  % levels in table of contents

% --- fancyhdr package for fancy headers ---
\usepackage{fancyhdr}
\fancyhf{} % sets both header and footer to nothing
\renewcommand{\headrulewidth}{0pt}
\fancyfoot[LE,RO]{\thepage}
% Ensure copyright on titlepage (article style) and chapter pages (book style)
\fancypagestyle{plain}{
  \fancyhf{}
  \fancyfoot[C]{{\footnotesize \copyright\ 2019, Ahmed Ammar. Released under CC Attribution 4.0 license}}
%  \renewcommand{\footrulewidth}{0mm}
  \renewcommand{\headrulewidth}{0mm}
}
% Ensure copyright on titlepages with \thispagestyle{empty}
\fancypagestyle{empty}{
  \fancyhf{}
  \fancyfoot[C]{{\footnotesize \copyright\ 2019, Ahmed Ammar. Released under CC Attribution 4.0 license}}
  \renewcommand{\footrulewidth}{0mm}
  \renewcommand{\headrulewidth}{0mm}
}

\pagestyle{fancy}


\usepackage{framed,wrapfig}

% --- begin definitions of admonition environments ---

% Admonition style "grayicon" has colored background, no frame, and an icon
% Admon "notice"
\definecolor{grayicon_notice_background}{rgb}{0.91,0.91,0.91}
% \fboxsep sets the space between the text and the box
\newenvironment{noticeshaded}
{\def\FrameCommand{\fboxsep=3mm\colorbox{grayicon_notice_background}}
 \MakeFramed {\advance\hsize-\width \FrameRestore}}{\endMakeFramed}

\newenvironment{notice_grayiconadmon}[1][Notice]{
\begin{noticeshaded}
\noindent
\begin{wrapfigure}{l}{0.07\textwidth}
\vspace{-13pt}
\includegraphics[width=0.07\textwidth]{latex_figs/small_gray_notice}
\end{wrapfigure} \textbf{#1}\par
\nobreak\noindent\ignorespaces
}
{
\end{noticeshaded}
}

% Admonition style "grayicon" has colored background, no frame, and an icon
% Admon "summary"
\definecolor{grayicon_summary_background}{rgb}{0.91,0.91,0.91}
% \fboxsep sets the space between the text and the box
\newenvironment{summaryshaded}
{\def\FrameCommand{\fboxsep=3mm\colorbox{grayicon_summary_background}}
 \MakeFramed {\advance\hsize-\width \FrameRestore}}{\endMakeFramed}

\newenvironment{summary_grayiconadmon}[1][Summary]{
\begin{summaryshaded}
\noindent
\begin{wrapfigure}{l}{0.07\textwidth}
\vspace{-13pt}
\includegraphics[width=0.07\textwidth]{latex_figs/small_gray_summary}
\end{wrapfigure} \textbf{#1}\par
\nobreak\noindent\ignorespaces
}
{
\end{summaryshaded}
}

% Admonition style "grayicon" has colored background, no frame, and an icon
% Admon "warning"
\definecolor{grayicon_warning_background}{rgb}{0.91,0.91,0.91}
% \fboxsep sets the space between the text and the box
\newenvironment{warningshaded}
{\def\FrameCommand{\fboxsep=3mm\colorbox{grayicon_warning_background}}
 \MakeFramed {\advance\hsize-\width \FrameRestore}}{\endMakeFramed}

\newenvironment{warning_grayiconadmon}[1][Warning]{
\begin{warningshaded}
\noindent
\begin{wrapfigure}{l}{0.07\textwidth}
\vspace{-13pt}
\includegraphics[width=0.07\textwidth]{latex_figs/small_gray_warning}
\end{wrapfigure} \textbf{#1}\par
\nobreak\noindent\ignorespaces
}
{
\end{warningshaded}
}

% Admonition style "grayicon" has colored background, no frame, and an icon
% Admon "question"
\definecolor{grayicon_question_background}{rgb}{0.91,0.91,0.91}
% \fboxsep sets the space between the text and the box
\newenvironment{questionshaded}
{\def\FrameCommand{\fboxsep=3mm\colorbox{grayicon_question_background}}
 \MakeFramed {\advance\hsize-\width \FrameRestore}}{\endMakeFramed}

\newenvironment{question_grayiconadmon}[1][Question]{
\begin{questionshaded}
\noindent
\begin{wrapfigure}{l}{0.07\textwidth}
\vspace{-13pt}
\includegraphics[width=0.07\textwidth]{latex_figs/small_gray_question2}
\end{wrapfigure} \textbf{#1}\par
\nobreak\noindent\ignorespaces
}
{
\end{questionshaded}
}

% Admonition style "grayicon" has colored background, no frame, and an icon
% Admon "block"
\definecolor{grayicon_block_background}{rgb}{0.91,0.91,0.91}
% \fboxsep sets the space between the text and the box
\newenvironment{blockshaded}
{\def\FrameCommand{\fboxsep=3mm\colorbox{grayicon_block_background}}
 \MakeFramed {\advance\hsize-\width \FrameRestore}}{\endMakeFramed}

\newenvironment{block_grayiconadmon}[1][Block]{
\begin{blockshaded}
\noindent
 \textbf{#1}\par
\nobreak\noindent\ignorespaces
}
{
\end{blockshaded}
}

% --- end of definitions of admonition environments ---

% prevent orhpans and widows
\clubpenalty = 10000
\widowpenalty = 10000

% --- end of standard preamble for documents ---


% insert custom LaTeX commands...

\raggedbottom
\makeindex
\usepackage[totoc]{idxlayout}   % for index in the toc
\usepackage[nottoc]{tocbibind}  % for references/bibliography in the toc

%-------------------- end preamble ----------------------

\begin{document}

% matching end for #ifdef PREAMBLE
% #endif

\newcommand{\exercisesection}[1]{\subsection*{#1}}


% ------------------- main content ----------------------



% ----------------- title -------------------------

\title{Apprenez à développer des applications de bureau avec Python and Qt}

% ----------------- author(s) -------------------------

\author{Ahmed Ammar\footnote{Email: \texttt{ahmed.ammar@fst.utm.tn}. Faculté des Sciences de Tunis, Université de Tunis El Manar.}}

% ----------------- end author(s) -------------------------

\date{Feb 15, 2019}
\maketitle

\tableofcontents


\vspace{1cm} % after toc




% !split
\section{Introduction}
Dans ce chapitre, nous allons passer à la création d’applications graphiques. De telles applications permettent de modifier l'apparence d'un programme en utilisant des \textbf{éléments de contrôle} tels que des \textbf{widgets}, des \textbf{boutons}, des \textbf{curseurs}, etc.

\textbf{PyQt5} (\href{{https://www.riverbankcomputing.com/software/pyqt/intro}}{\nolinkurl{https://www.riverbankcomputing.com/software/pyqt/intro}}) est une boîte à outils de widgets d' interface graphique (en anglais \textbf{GUI} pour \emph{Graphical User Interface}) combinant le langage de programmation Python et le logiciel \textbf{Qt5} (\href{{https://www.qt.io}}{\nolinkurl{https://www.qt.io}}). PyQt5 est une bibliothèque d'interface graphique populaire qui présente de nombreux avantages par rapport aux autres bibliothèques d'interface graphique telles que \emph{Tkinter} et \emph{wx}. Ceux-ci inclus:

\begin{itemize}
\item Bibliothèque d'interface graphique multi-plateforme (Windows, MacOS, Linux).

\item Bonne performance.

\item Prise en charge des styles personnalisés.

\item Plus de bibliothèques pour la conception d'interface graphique complexe.

\item Facilité d'utilisation.
\end{itemize}

\noindent
\section{Premières étapes pour la création d'une application graphique à l'aide de PyQt5}

\subsection{Importation de PyQt5 et création d'une "fenêtre PyQt5"}

Passons directement à l'action et apprenons à créer une fenêtre simple avec PyQt5. Premièrement, nous devrons importer certains modules essentiels à l’exécution d’une interface graphique avec PyQt5. Nous commençons par \textbf{importer} quelques sous-modules de PyQt5.

\begin{cod}{cbg_gray}\begin{minted}[fontsize=\fontsize{9pt}{9pt},linenos=false,mathescape,baselinestretch=1.0,fontfamily=tt,xleftmargin=2mm]{python}
from PyQt5.QtWidgets import QApplication, QWidget
\end{minted}
\end{cod}
\noindent

\begin{notice_grayiconadmon}[Notice]
\texttt{QtWidgets} est l'un des nombreux composants de PyQt5. Certains des plus couramment utilisés sont énumérés ci-dessous:
\begin{itemize}
\item \texttt{QtWidgets}: Contient des classes qui fournissent un ensemble d'éléments pour créer une interface graphique classique de type bureau.

\item \texttt{QtCore}: Ce module contient les classes principales, y compris la boucle d'événement et les mécanismes \texttt{signal} et \texttt{slot} de Qt.

\item \texttt{QtGui}: Celui-ci contient des classes pour l’intégration du système de fenêtrage, la gestion des événements, les graphiques 2D, les images de base, les polices et le texte.

\item \texttt{QtDesigner}: Ce module contient des classes permettant l'extension de \emph{Qt Designer} à l'aide de PyQt5.

\item \texttt{uic}: Ce module contient des classes permettant de gérer les fichiers \textbf{.ui} créés par Qt Designer décrivant l’ensemble ou une partie d’une interface graphique.

\item etc.
\end{itemize}

\noindent
\end{notice_grayiconadmon} % title: Notice


Ensuite, nous importons le module \texttt{sys} car nous voulons accéder aux arguments de ligne de commande. Ceux-ci sont contenus dans la liste `sys.argv

\begin{cod}{cbg_gray}\begin{minted}[fontsize=\fontsize{9pt}{9pt},linenos=false,mathescape,baselinestretch=1.0,fontfamily=tt,xleftmargin=2mm]{python}
import sys
app = QApplication(sys.argv)
\end{minted}
\end{cod}
\noindent

En utilisant les arguments de ligne de commande contenus dans \texttt{sys.argv}, nous devrions créer un objet \texttt{QApplication}. Dans l'exemple ci-dessus, nous avons enregistré cet objet sous le nom \textbf{app}. Ceci est l'objet \emph{exécuté} lorsque nous exécutons à l'invite de commande la commande \texttt{python filename.py} où \textbf{filename.py} est le fichier dans lequel les instructions sont stockées. Cet objet contiendra tous les éléments de l'interface graphique et leurs méthodes. Par conséquent, lorsque cet objet est créé et \emph{exécuté}, nous avons accès à toutes les interfaces graphiques utilisées dans le programme. Les éléments de l'interface graphique sont hérités de la classe \texttt{QWidgets}. Nous obtenons donc un objet \texttt{QWidget} comme indiqué ci-dessous.
\begin{cod}{cbg_gray}\begin{minted}[fontsize=\fontsize{9pt}{9pt},linenos=false,mathescape,baselinestretch=1.0,fontfamily=tt,xleftmargin=2mm]{python}
window = QWidget()
\end{minted}
\end{cod}
\noindent
En utilisant les nombreuses méthodes \emph{set}, nous pouvons définir les valeurs des différents attributs de cet objet \emph{window}. Nous allons d’abord définir les dimensions de la fenêtre en utilisant la méthode \texttt{setGeometry}. Il prend quatre paramètres: \textbf{x} coordonnée du point le plus à gauche de l'objet (\emph{window}), \textbf{y} coordonnée du point le plus haut, \textbf{la largeur} et \textbf{la hauteur} de la fenêtre, dans cet ordre. Toutes ces valeurs doivent être des entiers référencés par rapport aux coordonnées d'écran natives. Par exemple, nous pourrions écrire:

\begin{cod}{cbg_gray}\begin{minted}[fontsize=\fontsize{9pt}{9pt},linenos=false,mathescape,baselinestretch=1.0,fontfamily=tt,xleftmargin=2mm]{python}
window.setGeometry(400, 100, 300, 200)
\end{minted}
\end{cod}
\noindent
Cela définirait la fenêtre avec une largeur de 300 pixels et une hauteur de 200 pixels à 400 pixels du côté gauche et à 100 pixels du haut de l'écran natif (écran de votre ordinateur).

Nous pouvons définir un titre pour la fenêtre. Par exemple:

\begin{cod}{cbg_gray}\begin{minted}[fontsize=\fontsize{9pt}{9pt},linenos=false,mathescape,baselinestretch=1.0,fontfamily=tt,xleftmargin=2mm]{python}
window.setWindowTitle('My first app')
\end{minted}
\end{cod}
\noindent
afficherait \emph{My first app} dans la barre de titre de la fenêtre.

Une fois que nous sommes satisfaits de l'interface graphique que nous avons construite, nous appelons:
\begin{cod}{cbg_gray}\begin{minted}[fontsize=\fontsize{9pt}{9pt},linenos=false,mathescape,baselinestretch=1.0,fontfamily=tt,xleftmargin=2mm]{python}
window.show()
\end{minted}
\end{cod}
\noindent
pour afficher l'objet graphique dans \emph{l'application} que nous avons créée. Cependant, l'application n'est toujours pas exécutée. Pour exécuter l'application, nous exécutons la commande:
\begin{cod}{cbg_gray}\begin{minted}[fontsize=\fontsize{9pt}{9pt},linenos=false,mathescape,baselinestretch=1.0,fontfamily=tt,xleftmargin=2mm]{python}
app.exec()
\end{minted}
\end{cod}
\noindent
La méthode \texttt{.show ()} est une méthode QT qui ouvre la fenêtre à l'écran pour l'utilisateur.

Enfin, pour assurer une fermeture agréable et propre de l'application lorsque nous quittons:
\begin{cod}{cbg_gray}\begin{minted}[fontsize=\fontsize{9pt}{9pt},linenos=false,mathescape,baselinestretch=1.0,fontfamily=tt,xleftmargin=2mm]{python}
sys.exit(app.exec_())
\end{minted}
\end{cod}
\noindent
Le code complet de cette application est donné ci-dessous et le résultat est illustré à la figure 1. Le code peut être stocké dans un fichier, par exemple \textbf{myfirstapp.py}, puis lorsque vous exécutez la commande \texttt{python MyFirstApp.py} à l'invite de commande (ou dans Spyder), vous obtenez la fenêtre de sortie affichée.

\begin{pro}{cbg_gray}{bar_gray}\begin{minted}[fontsize=\fontsize{9pt}{9pt},linenos=false,mathescape,baselinestretch=1.0,fontfamily=tt,xleftmargin=2mm]{python}
# NOM DU FICHIER: myfirstapp.py
#% IMPORTATION
from PyQt5.QtWidgets import QApplication, QWidget
import sys
# créer un objet QApplication
app = QApplication(sys.argv)
# appeler la classe QWidgets
window = QWidget()
# définir les dimensions de la fenêtre
window.setGeometry(400, 100, 300, 200)
# définir un titre pour la fenêtre
window.setWindowTitle('My first app')
# afficher l'objet graphique dans l'application
window.show()
# fermer l'application
sys.exit(app.exec_())
\end{minted}
\end{pro}
\noindent

% ------------------- end of main content ---------------

% #ifdef PREAMBLE
\end{document}
% #endif

